\documentclass[10pt,a4paper]{amsart}
\usepackage[margin=19mm]{geometry}
\usepackage{amsmath,amssymb,mathtools}
\usepackage[scr=rsfs,cal=euler]{mathalfa}
\usepackage{xfrac}
\usepackage[unicode,hidelinks,bookmarks=false]{hyperref}
\newcommand{\ie}{i.e.\ }
\newcommand{\cf}{cf.\ }
\newcommand{\resp}{respectively\ }
\newcommand{\Subsection}{\subsection}
\providecommand{\nobreakdash}{\nobreak\hbox{-}}
\setlength{\emergencystretch}{2em}
\multlinegap0pt
\usepackage{cleveref}
\newcommand{\R}{\mathbb R}
\newcommand{\calC}{\mathcal C}
\newcommand{\calF}{\mathcal F}
\newcommand{\calT}{\mathcal T}
\newcommand{\dd}{\,d}
\theoremstyle{plain}
\newtheorem{theorem}{Theorem}[section]
\newtheorem{lemma}[theorem]{Lemma}
\newtheorem{proposition}[theorem]{Proposition}
\newtheorem{corollary}[theorem]{Corollary}
\theoremstyle{definition}
\newtheorem{definition}[theorem]{Definition}
\newtheorem{assumption}[theorem]{Assumption}
\theoremstyle{remark}
\newtheorem{remark}[theorem]{Remark}
\title[Harnack-type estimates for nonlocal double phase functionals]{Harnack-type estimates for nonlocal double phase functionals}
\author{Yuzhou Fang, Chao Zhang, Yuzhou Fang and Chao Zhang$^*$}
\date{Source: arXiv:2609.27957v1 (CC BY 4.0)}
\begin{document}
\maketitle

\noindent\emph{Source and licence.} The delivered work is the article shown in the title above, version arXiv:2609.27957v1 (CC BY 4.0), distributed under the Creative Commons Attribution 4.0 International licence, as recorded in the arXiv licence metadata for this exact version and shown on the abstract page saved with this item. The full licence notice, which carries the licence URL and the address of that abstract page, is the bundled file \texttt{licenses/arxiv-2609.27957-CC-BY-4.0.txt}.
\par\smallskip

\noindent\textit{Editorial scope.} Counted unit: the article's Corollary (source label cor:harnack-single-tail) (source0.tex lines 384--391, source label \texttt{cor:harnack-single-tail}): ``Under the assumptions of , there exists a constant (C 1 ), depending only on ( D u(B 2R (x 0)) ), such that B R 4 (x 0) u C B R 4 (x 0) u (h 2R ) -1 ( R 4,2R ( u ;x 0) ) .''. It is delivered together with the article's complete original proof (source0.tex lines 1252--1360, one proof environment printed later in the article, detached from the statement by the article's own arrangement, 109 lines), copied line for line from the article's own text. Required context retained uncounted: main (lines 136--140); eq:exponents (lines 154--156); eq:sobolev-range (lines 180--186); thm:harnack-tail (lines 359--380). Editorial formatting only: an AMS \texttt{amsart} preamble that restates the article's own macro definitions and its own theorem declarations, and the theorem environments the retained lines use; the packages cleveref are loaded to match the article's own setup; the article's own class file is neither shipped nor input; the retained environments are renumbered sequentially in order of appearance, whereas the article prints them with its own numbering; the article's equation numbering is not reproduced and every cross-reference inside the retained text resolves to the wrapper's numbering. No citation occurs inside any retained line, so the wrapper carries no bibliography; All retained bodies are exact copies of the bundled \texttt{sources/211533/source0.tex}, so no omission receipt applies. No asset-inclusion command occurs inside any retained span and the package ships no image asset.


\section*{Retained context: main (source0.tex lines 136--140)}

\begin{equation}\label{main}
    \calF(u;\Omega)
    :=\iint_{\calC_\Omega}\left(|u(x)-u(y)|^p K_{sp}(x,y)+ a(x,y)|u(x)-u(y)|^q K_{tq}(x,y)
    \right)\dd x\!\dd y ,
\end{equation}


\section*{Retained context: eq:exponents (source0.tex lines 154--156)}

\begin{equation}\label{eq:exponents}
    0<s,t<1,\qquad 1<p\le q<\infty.
\end{equation}


\section*{Retained context: eq:sobolev-range (source0.tex lines 180--186)}

\begin{equation}\label{eq:sobolev-range}
    q\le\frac{np}{n-sp}
    \quad\text{if }sp<n,
    \qquad
    p\le q<\infty
    \quad\text{if }sp\ge n.
\end{equation}


\section*{Retained context: thm:harnack-tail (source0.tex lines 359--380)}

\begin{theorem}[Harnack-type inequality]
\label{thm:harnack-tail}
Assume that \eqref{eq:exponents}--\eqref{eq:sobolev-range} hold and
\begin{align*}
\begin{cases}
\alpha<tq \leq sp+\alpha,  \quad &\text{for  } s\le t,\\[2mm]
tq\le sp,        \quad&\text{for  } s>t.
\end{cases}
\end{align*}
Let \(u\in\mathcal{A}(\Omega)\cap L^{p-1}_{sp}(\R^n)
\cap L^{q-1}_{a,tq}(\Omega,\R^n)\) be a local minimizer
of~\eqref{main} in \(\Omega\). Let \(0<R\le1\), suppose that
\(B_{2R}(x_0)\Subset\Omega\) and \(u\ge0\) a.e.\ in
\(B_{2R}(x_0)\).
Then there exists a constant \(C\ge1\), depending only on
\(\mathfrak D_u(B_{2R}(x_0))\), such that
\begin{equation}\label{eq:harnack-tail}
    \sup_{B_{R/4}(x_0)}u
    \le
    C\left[\inf_{B_{R/4}(x_0)}u+(h_{R/2})^{-1}\bigl(\calT_{R/4,R/2}(u_+;x_0)\bigr)+(h_{2R})^{-1}\bigl(\calT_{2R}(u_-;x_0)\bigr)\right].
\end{equation}
\end{theorem}


\section{The counted unit: Corollary (source label cor:harnack-single-tail) and its complete original proof}

\begin{corollary}\label{cor:harnack-single-tail}
Under the assumptions of \Cref{thm:harnack-tail}, there exists a constant
\(C\ge1\), depending only on \(\mathfrak D_u(B_{2R}(x_0))\), such that
\begin{equation}\label{eq:harnack-single-tail}
    \sup_{B_{R/4}(x_0)}u\le
    C\left[\inf_{B_{R/4}(x_0)}u+(h_{2R})^{-1}\bigl(\calT_{R/4,2R}(|u|;x_0)\bigr)\right].
\end{equation}
\end{corollary}

\begin{proof}[\textbf{Proof of \Cref{cor:harnack-single-tail}}]
For brevity, set
\[
    A:=\calT_{R/4,R/2}(u_+;x_0),
    \qquad
    B:=\calT_{2R}(u_-;x_0),
    \qquad
    T:=\calT_{R/4,2R}(|u|;x_0).
\]
The inclusion \(B_{R/2}(x_0)\times B_{R/2}(x_0)\subset B_{2R}(x_0)\times B_{2R}(x_0)\) implies \(a_{R/2}^-\ge a_{2R}^-\).  Thus, for every $\tau\ge0$,
\[
\begin{aligned}
    h_{R/2}(\tau)
    &=\frac{\tau^{p-1}}{(R/2)^{sp}}
      +a_{R/2}^-\frac{\tau^{q-1}}{(R/2)^{tq}}\\
    &\ge
      \frac{\tau^{p-1}}{(2R)^{sp}}
      +a_{2R}^-\frac{\tau^{q-1}}{(2R)^{tq}}
      =h_{2R}(\tau).
\end{aligned}
\]
Since both functions are strictly increasing bijections of
\([0,\infty)\), their inverses satisfy
\[
    (h_{R/2})^{-1}(S)\le (h_{2R})^{-1}(S)
    \qquad\text{for every }S\ge0.
\]
The monotonicity of \((h_{2R})^{-1}\) therefore yields
\begin{equation}\label{eq:inverse-tail-comparison}
\begin{aligned}
    &\quad(h_{R/2})^{-1}(A)+(h_{2R})^{-1}(B)\\
    &\le
    (h_{2R})^{-1}(A)+(h_{2R})^{-1}(B)
    \le 2(h_{2R})^{-1}(A+B).
\end{aligned}
\end{equation}

The sign assumption in \Cref{thm:harnack-tail} gives \(u_-=0\) a.e.\ in
\(B_{2R}(x_0)\). Hence enlarging the exterior integration region from
\(\R^n\setminus B_{2R}(x_0)\) to \(\R^n\setminus B_{R/4}(x_0)\) adds only a region on which \(u_-=0\), so
\[
    \calT_{2R}(u_-;x_0)
    =\calT_{R/4,2R}(u_-;x_0),
\]
where the interior reference ball is \(B_{2R}(x_0)\) on both sides. Moreover, enlarging the ball over which the weighted \(q\)-tail supremum
is taken gives
\[
    \calT_{R/4,R/2}(u_+;x_0)
    \le \calT_{R/4,2R}(u_+;x_0).
\]

To verify the resulting signed-tail comparison directly, define
\[
\begin{aligned}
    P(v)
    &:={}
      \int_{\R^n\setminus B_{R/4}(x_0)}
      \frac{|v(y)|^{p-1}}{|x_0-y|^{n+sp}}\dd y,\\
    Q(v;x)
    &:={}
      \int_{\R^n\setminus B_{R/4}(x_0)}
      a(x,y)\frac{|v(y)|^{q-1}}{|x_0-y|^{n+tq}}\dd y,
      \qquad x\in B_{2R}(x_0).
\end{aligned}
\]
The preceding two observations imply
\[
    A+B
    \le P(u_+)+P(u_-)+
       \sup_{x\in B_{2R}(x_0)}Q(u_+;x)
       +\sup_{x\in B_{2R}(x_0)}Q(u_-;x).
\]
Since \(|u|^{m-1}=u_+^{m-1}+u_-^{m-1}\) for \(m=p,q\), we have
\(P(u_+)+P(u_-)=P(|u|)\) and
\(Q(u_+;x)+Q(u_-;x)=Q(|u|;x)\). It follows that
\begin{equation}\label{eq:signed-tail-comparison}
\begin{aligned}
    A+B
    &\le P(|u|)+2\sup_{x\in B_{2R}(x_0)}Q(|u|;x)\\
    &\le2\calT_{R/4,2R}(|u|;x_0)=2T.
\end{aligned}
\end{equation}

Since \(q\ge p\), for every \(\lambda\ge1\) and \(\tau\ge0\),
\[
\begin{aligned}
    h_{2R}(\lambda\tau)=\lambda^{p-1}\frac{\tau^{p-1}}{(2R)^{sp}}
      +a_{2R}^-\lambda^{q-1}\frac{\tau^{q-1}}{(2R)^{tq}}
    \ge \lambda^{p-1}h_{2R}(\tau).
\end{aligned}
\]
Taking \(\tau=(h_{2R})^{-1}(S)\) and
\(\lambda=2^{1/(p-1)}\), and then applying the increasing inverse, gives
\[
    (h_{2R})^{-1}(2S)
    \le 2^{1/(p-1)}(h_{2R})^{-1}(S)
    \qquad\text{for every }S\ge0.
\]
Combining this with \eqref{eq:inverse-tail-comparison} and \eqref{eq:signed-tail-comparison}, we obtain the explicit estimate
\[
\begin{aligned}
    &\quad (h_{R/2})^{-1}(A)+(h_{2R})^{-1}(B)\\
    &\le2(h_{2R})^{-1}(A+B)
     \le2(h_{2R})^{-1}(2T)
     \le2^{1+1/(p-1)}(h_{2R})^{-1}(T).
\end{aligned}
\]
Substitution of this bound into \eqref{eq:harnack-tail} proves \eqref{eq:harnack-single-tail}. %; the fixed factor \(2^{1+1/(p-1)}\), which depends only on the structural exponent \(p\), is absorbed into the constant \(C\).
\end{proof}

\end{document}
